\subsection{Earthing systems}
\label{sec:de-earthing}

DIN~VDE~0100-100 designates the earthing system by two
letters~[20]. The first gives the relation of the source to
earth: T, one point (normally the neutral) directly earthed; I, all live
parts isolated from earth or one point earthed through a high impedance. The
second gives the earthing of the exposed conductive parts of the
installation: N, connected to the earthed point of the source through
protective conductors; T, connected to an earth electrode independent of
the source earth. For TN systems, C denotes a combined protective earth and
neutral conductor (PEN) and S separate N and PE conductors.
Figure~\ref{fig:de-earthing} shows the five systems.

\begin{figure}[!ht]
  \centering
  % generated by north-american-practice/fig/connection_figures.py -- do not edit
\begin{tikzpicture}[x=0.22mm, y=-0.22mm, line cap=round, line join=round]
  \draw[box] (10.0,20.0) rectangle (44.0,104.0);
  \node[font=\scriptsize, text=atFgMuted, anchor=center, inner sep=0.5pt, rotate=90] at (27.0,62.0) {Dyn5};
  \node[font=\small\bfseries, text=atFg, anchor=west, inner sep=0.5pt] at (10.0,6.0) {TN-S};
  \draw[ph] (44.0,30.0) -- (300.0,30.0);
  \node[font=\scriptsize, text=atFgMuted, anchor=west, inner sep=0.5pt] at (305.0,30.0) {L1};
  \draw[ph] (44.0,44.0) -- (300.0,44.0);
  \node[font=\scriptsize, text=atFgMuted, anchor=west, inner sep=0.5pt] at (305.0,44.0) {L2};
  \draw[ph] (44.0,58.0) -- (300.0,58.0);
  \node[font=\scriptsize, text=atFgMuted, anchor=west, inner sep=0.5pt] at (305.0,58.0) {L3};
  \draw[box] (200.0,118.0) rectangle (278.0,152.0);
  \node[font=\scriptsize, text=atFgMuted, anchor=center, inner sep=0.5pt] at (239.0,135.0) {equipment};
  \draw[ph] (214.0,30.0) -- (214.0,118.0);
  \fill[atFg] (214.0,30.0) circle (3.0);
  \draw[nl] (44.0,76.0) -- (300.0,76.0);
  \draw[pey] (56.0,94.0) -- (300.0,94.0);
  \draw[peg] (56.0,94.0) -- (300.0,94.0);
  \draw[pey] (56.0,76.0) -- (56.0,94.0);
  \node[font=\scriptsize, text=atInfo, anchor=west, inner sep=0.5pt] at (305.0,76.0) {N};
  \node[font=\scriptsize, text=atSuccess, anchor=west, inner sep=0.5pt] at (305.0,94.0) {PE};
  \fill[atInfo] (56.0,76.0) circle (3.0);
  \draw[gnd] (56.0,76.0) -- (56.0,140.0);
  \draw[gnd] (56.0,140.0) -- (56.0,150.0);
  \draw[gnd] (49.0,150.0) -- (63.0,150.0);
  \draw[gnd] (51.5,154.0) -- (60.5,154.0);
  \draw[gnd] (54.0,158.0) -- (58.0,158.0);
  \draw[nl] (234.0,76.0) -- (234.0,118.0);
  \fill[atInfo] (234.0,76.0) circle (3.0);
  \draw[pey] (262.0,94.0) -- (262.0,118.0);
  \draw[peg] (262.0,94.0) -- (262.0,118.0);
  \draw[box] (370.0,20.0) rectangle (404.0,104.0);
  \node[font=\scriptsize, text=atFgMuted, anchor=center, inner sep=0.5pt, rotate=90] at (387.0,62.0) {Dyn5};
  \node[font=\small\bfseries, text=atFg, anchor=west, inner sep=0.5pt] at (370.0,6.0) {TN-C};
  \draw[ph] (404.0,30.0) -- (660.0,30.0);
  \node[font=\scriptsize, text=atFgMuted, anchor=west, inner sep=0.5pt] at (665.0,30.0) {L1};
  \draw[ph] (404.0,44.0) -- (660.0,44.0);
  \node[font=\scriptsize, text=atFgMuted, anchor=west, inner sep=0.5pt] at (665.0,44.0) {L2};
  \draw[ph] (404.0,58.0) -- (660.0,58.0);
  \node[font=\scriptsize, text=atFgMuted, anchor=west, inner sep=0.5pt] at (665.0,58.0) {L3};
  \draw[box] (560.0,118.0) rectangle (638.0,152.0);
  \node[font=\scriptsize, text=atFgMuted, anchor=center, inner sep=0.5pt] at (599.0,135.0) {equipment};
  \draw[ph] (574.0,30.0) -- (574.0,118.0);
  \fill[atFg] (574.0,30.0) circle (3.0);
  \draw[nl] (404.0,76.0) -- (660.0,76.0);
  \draw[peg, line width=0.8pt] (404.0,76.0) -- (660.0,76.0);
  \node[font=\scriptsize, text=atSuccess, anchor=west, inner sep=0.5pt] at (665.0,76.0) {PEN};
  \fill[atInfo] (416.0,76.0) circle (3.0);
  \draw[gnd] (416.0,76.0) -- (416.0,140.0);
  \draw[gnd] (416.0,140.0) -- (416.0,150.0);
  \draw[gnd] (409.0,150.0) -- (423.0,150.0);
  \draw[gnd] (411.5,154.0) -- (420.5,154.0);
  \draw[gnd] (414.0,158.0) -- (418.0,158.0);
  \draw[nl] (594.0,76.0) -- (594.0,118.0);
  \fill[atInfo] (594.0,76.0) circle (3.0);
  \draw[peg, line width=0.8pt] (622.0,76.0) -- (622.0,118.0);
  \fill[atInfo] (622.0,76.0) circle (3.0);
  \draw[box] (10.0,220.0) rectangle (44.0,304.0);
  \node[font=\scriptsize, text=atFgMuted, anchor=center, inner sep=0.5pt, rotate=90] at (27.0,262.0) {Dyn5};
  \node[font=\small\bfseries, text=atFg, anchor=west, inner sep=0.5pt] at (10.0,206.0) {TN-C-S};
  \draw[ph] (44.0,230.0) -- (300.0,230.0);
  \node[font=\scriptsize, text=atFgMuted, anchor=west, inner sep=0.5pt] at (305.0,230.0) {L1};
  \draw[ph] (44.0,244.0) -- (300.0,244.0);
  \node[font=\scriptsize, text=atFgMuted, anchor=west, inner sep=0.5pt] at (305.0,244.0) {L2};
  \draw[ph] (44.0,258.0) -- (300.0,258.0);
  \node[font=\scriptsize, text=atFgMuted, anchor=west, inner sep=0.5pt] at (305.0,258.0) {L3};
  \draw[box] (200.0,318.0) rectangle (278.0,352.0);
  \node[font=\scriptsize, text=atFgMuted, anchor=center, inner sep=0.5pt] at (239.0,335.0) {equipment};
  \draw[ph] (214.0,230.0) -- (214.0,318.0);
  \fill[atFg] (214.0,230.0) circle (3.0);
  \draw[nl] (44.0,276.0) -- (300.0,276.0);
  \draw[peg, line width=0.8pt] (44.0,276.0) -- (150.0,276.0);
  \draw[pey] (150.0,294.0) -- (300.0,294.0);
  \draw[peg] (150.0,294.0) -- (300.0,294.0);
  \draw[pey] (150.0,276.0) -- (150.0,294.0);
  \draw[peg] (150.0,276.0) -- (150.0,294.0);
  \fill[atFg] (150.0,276.0) circle (3.0);
  \node[font=\scriptsize, text=atFgMuted, anchor=center, inner sep=0.5pt] at (150.0,268.0) {split};
  \node[font=\scriptsize, text=atSuccess, anchor=center, inner sep=0.5pt] at (86.0,285.0) {PEN};
  \node[font=\scriptsize, text=atInfo, anchor=west, inner sep=0.5pt] at (305.0,276.0) {N};
  \node[font=\scriptsize, text=atSuccess, anchor=west, inner sep=0.5pt] at (305.0,294.0) {PE};
  \fill[atInfo] (56.0,276.0) circle (3.0);
  \draw[gnd] (56.0,276.0) -- (56.0,340.0);
  \draw[gnd] (56.0,340.0) -- (56.0,350.0);
  \draw[gnd] (49.0,350.0) -- (63.0,350.0);
  \draw[gnd] (51.5,354.0) -- (60.5,354.0);
  \draw[gnd] (54.0,358.0) -- (58.0,358.0);
  \draw[nl] (234.0,276.0) -- (234.0,318.0);
  \fill[atInfo] (234.0,276.0) circle (3.0);
  \draw[gnd] (150.0,294.0) -- (150.0,340.0);
  \draw[gnd] (150.0,340.0) -- (150.0,350.0);
  \draw[gnd] (143.0,350.0) -- (157.0,350.0);
  \draw[gnd] (145.5,354.0) -- (154.5,354.0);
  \draw[gnd] (148.0,358.0) -- (152.0,358.0);
  \draw[pey] (262.0,294.0) -- (262.0,318.0);
  \draw[peg] (262.0,294.0) -- (262.0,318.0);
  \draw[box] (370.0,220.0) rectangle (404.0,304.0);
  \node[font=\scriptsize, text=atFgMuted, anchor=center, inner sep=0.5pt, rotate=90] at (387.0,262.0) {Dyn5};
  \node[font=\small\bfseries, text=atFg, anchor=west, inner sep=0.5pt] at (370.0,206.0) {TT};
  \draw[ph] (404.0,230.0) -- (660.0,230.0);
  \node[font=\scriptsize, text=atFgMuted, anchor=west, inner sep=0.5pt] at (665.0,230.0) {L1};
  \draw[ph] (404.0,244.0) -- (660.0,244.0);
  \node[font=\scriptsize, text=atFgMuted, anchor=west, inner sep=0.5pt] at (665.0,244.0) {L2};
  \draw[ph] (404.0,258.0) -- (660.0,258.0);
  \node[font=\scriptsize, text=atFgMuted, anchor=west, inner sep=0.5pt] at (665.0,258.0) {L3};
  \draw[box] (560.0,318.0) rectangle (638.0,352.0);
  \node[font=\scriptsize, text=atFgMuted, anchor=center, inner sep=0.5pt] at (599.0,335.0) {equipment};
  \draw[ph] (574.0,230.0) -- (574.0,318.0);
  \fill[atFg] (574.0,230.0) circle (3.0);
  \draw[nl] (404.0,276.0) -- (660.0,276.0);
  \node[font=\scriptsize, text=atInfo, anchor=west, inner sep=0.5pt] at (665.0,276.0) {N};
  \fill[atInfo] (416.0,276.0) circle (3.0);
  \draw[gnd] (416.0,276.0) -- (416.0,340.0);
  \draw[gnd] (416.0,340.0) -- (416.0,350.0);
  \draw[gnd] (409.0,350.0) -- (423.0,350.0);
  \draw[gnd] (411.5,354.0) -- (420.5,354.0);
  \draw[gnd] (414.0,358.0) -- (418.0,358.0);
  \draw[nl] (594.0,276.0) -- (594.0,318.0);
  \fill[atInfo] (594.0,276.0) circle (3.0);
  \draw[pey] (622.0,352.0) -- (622.0,360.0) -- (660.0,360.0);
  \draw[peg] (622.0,352.0) -- (622.0,360.0) -- (660.0,360.0);
  \draw[gnd] (660.0,360.0) -- (660.0,370.0);
  \draw[gnd] (653.0,370.0) -- (667.0,370.0);
  \draw[gnd] (655.5,374.0) -- (664.5,374.0);
  \draw[gnd] (658.0,378.0) -- (662.0,378.0);
  \node[font=\scriptsize, text=atSuccess, anchor=center, inner sep=0.5pt] at (652.0,350.0) {local};
  \draw[box] (10.0,420.0) rectangle (44.0,504.0);
  \node[font=\scriptsize, text=atFgMuted, anchor=center, inner sep=0.5pt, rotate=90] at (27.0,462.0) {Dyn5};
  \node[font=\small\bfseries, text=atFg, anchor=west, inner sep=0.5pt] at (10.0,406.0) {IT};
  \draw[ph] (44.0,430.0) -- (300.0,430.0);
  \node[font=\scriptsize, text=atFgMuted, anchor=west, inner sep=0.5pt] at (305.0,430.0) {L1};
  \draw[ph] (44.0,444.0) -- (300.0,444.0);
  \node[font=\scriptsize, text=atFgMuted, anchor=west, inner sep=0.5pt] at (305.0,444.0) {L2};
  \draw[ph] (44.0,458.0) -- (300.0,458.0);
  \node[font=\scriptsize, text=atFgMuted, anchor=west, inner sep=0.5pt] at (305.0,458.0) {L3};
  \draw[box] (200.0,518.0) rectangle (278.0,552.0);
  \node[font=\scriptsize, text=atFgMuted, anchor=center, inner sep=0.5pt] at (239.0,535.0) {equipment};
  \draw[ph] (214.0,430.0) -- (214.0,518.0);
  \fill[atFg] (214.0,430.0) circle (3.0);
  \draw[box] (47.0,482.0) rectangle (65.0,504.0);
  \node[font=\scriptsize, text=atFgMuted, anchor=center, inner sep=0.5pt] at (56.0,493.0) {Z};
  \draw[gnd] (56.0,458.0) -- (56.0,482.0);
  \draw[gnd] (56.0,504.0) -- (56.0,540.0);
  \draw[gnd] (56.0,540.0) -- (56.0,550.0);
  \draw[gnd] (49.0,550.0) -- (63.0,550.0);
  \draw[gnd] (51.5,554.0) -- (60.5,554.0);
  \draw[gnd] (54.0,558.0) -- (58.0,558.0);
  \fill[atFg] (56.0,458.0) circle (3.0);
  \node[font=\scriptsize, text=atFgMuted, anchor=west, inner sep=0.5pt] at (70.0,493.0) {isolated or high impedance};
  \draw[ph] (234.0,444.0) -- (234.0,518.0);
  \fill[atFg] (234.0,444.0) circle (3.0);
  \draw[pey] (262.0,552.0) -- (262.0,560.0) -- (300.0,560.0);
  \draw[peg] (262.0,552.0) -- (262.0,560.0) -- (300.0,560.0);
  \draw[gnd] (300.0,560.0) -- (300.0,570.0);
  \draw[gnd] (293.0,570.0) -- (307.0,570.0);
  \draw[gnd] (295.5,574.0) -- (304.5,574.0);
  \draw[gnd] (298.0,578.0) -- (302.0,578.0);
  \node[font=\scriptsize, text=atSuccess, anchor=center, inner sep=0.5pt] at (292.0,550.0) {local};
\end{tikzpicture}

  \caption{Earthing systems: source on the left, equipment with its exposed
  conductive part on the right. Blue: neutral; green-yellow: protective
  earth; green dashed on blue: PEN.}
  \label{fig:de-earthing}
\end{figure}

\subsubsection{TN-S}
N and PE are separate from the transformer onwards and are connected only
at the earthed neutral point of the source. No operating current flows on
the PE, which is the reason TN-S is required or recommended for
installations with extensive information technology
equipment~[21].

\subsubsection{TN-C}
A combined PEN conductor serves as neutral and protective conductor. A PEN
conductor requires a cross-section of at least 10\,mm$^2$ copper or
16\,mm$^2$ aluminium in fixed installations~[22]; inside
buildings TN-C is limited to such conductors.

\subsubsection{TN-C-S}
The PEN conductor of the network is split into N and PE at a point of the
installation, normally the service entrance, where the PE is earthed again;
beyond the split, N and PE are not reconnected. Public low-voltage networks
in Germany are operated as TN-C systems up to the customer's service
entrance and continue as TN-S inside the installation, which makes the
combination TN-C-S~[13].

\subsubsection{TT}
The source neutral is earthed at the transformer, and the exposed
conductive parts are connected to a local earth electrode of the
installation. The earth-fault current returns through the soil, and
protection by automatic disconnection uses residual current devices with
\begin{equation}
  R_{a} \cdot I_{\Delta n} \le 50\,\mathrm{V}
  \label{eq:de-tt}
\end{equation}
where $R_{a}$ is the resistance of the local earth electrode and the
protective conductor, and $I_{\Delta n}$ the rated residual operating
current of the RCD~[23]. Some rural networks are operated as TT
systems.

\subsubsection{IT}
The source is isolated from earth or earthed through a high impedance, and
the exposed conductive parts are earthed locally. A first fault to earth
draws only the capacitive current of the system and does not require
disconnection; it is signalled by an insulation monitoring device. A
second fault on another phase is a short circuit and is disconnected as in
a TN or TT system~[23].
\FloatBarrier
